\section{Background and restatement of problem}
\label{sec:restate}

We start by introducing some of the background
material that leads to some reformulations of the
main theorem we wish to prove. We first introduce
the notions of ``constraints'' and ``(single-orbit)
characterizations,'' which leads to a first reformulation
of our main theorem (see \expref{Theorem}{thm:main}).
We then give some sufficient
conditions to recognize such characterizations,
and this leads to a second reformulation of our
main theorem (see \expref{Theorem}{thm:main2}).


\subsection{Single-orbit characterizations}

In this section we use the fact that Reed-Muller codes
form a ``linear, affine-invariant property.'' We recall
these notions first.
Given a finite field $\F_q$ a
property is a set of functions $\cF$ mapping $\F_{q}^n$ to $\F_q$.
The property is said to
be
\emph{linear} if
it is an $\F_q$-vector space, \ie, $\forall f,g \in \calf$
and $\alpha \in \F_q$ we have $\alpha f + g \in \calf$.
The property is said to be \emph{affine-invariant} if it is invariant under
affine-transformations of the domain, \ie,
$\forall f \in \calf$ it is the case that $f \circ T$ is also
in $\calf$ for every affine-transformation $T: \F_q^n \to \F_q^n$ given by $T(x) = A \cdot x + \beta$ for $A \in \F_q^{n \times n}$, $\beta \in \F_q^n$.\footnote{We note that as in~\cite{KS08} we do not require $A$ to be
non-singular. Thus the affine-transformations we consider are not
necessarily permutations from $\F_q^n$ to $\F_q^n$.}
It can be easily verified that $\mrm[n,d,q]$ is linear and affine-invariant
for every $n, d, q$.

The main tool used so far for constructing testers for affine-invariant linear properties is a
 structural theorem which shows that every linear affine-invariant property that is $k$-single-orbit characterizable
 is also $k$-locally testable.
 In order to describe the notion of single-orbit characterization we start with a couple of definitions.


\begin{definition}[$k$-constraint, $k$-characterization]
A \emph{$k$-constraint} $C=\bigl(\overline \alpha, \bset{\overline \lambda_i}_{i=1}^{r}\bigr)$
on $\{\F_{q}^n \to \F_q\}$ is
given by a vector $\overline \alpha = (\alpha_1,\ldots,\alpha_k) \in (\F_{q}^n)^k$
together with $r$ vectors
 $\overline \lambda_i = (\lambda_{i,1}, \ldots, \lambda_{i,k}) \in \F_q^k$ for
$1 \leq i \leq r$.
We say that the constraint $C$ \emph{accepts} a function $f: \F_{q}^n \to \F_{q}$ if
$\sum_{j=1}^k \lambda_{i,j} f(\alpha_j) = 0$ for all $1 \leq i \leq r$. Otherwise we say that $C$ \emph{rejects} $f$.

Let $\mathcal{F} \subseteq \set{\F_{q}^n \to \F_q}$ be a linear property.
A \emph{$k$-characterization} of $\calf$ is a collection of $k$-constraints
$C_1, \ldots, C_m$ on $\{\F_q^n \to \F_q\}$ such that $f \in \calf$ if and only if $C_j$
accepts $f$, for every $j \in \{1,\ldots,m\}$.
\end{definition}

It is well-known~\cite{BHR05} that every $k$-locally testable linear property must have a
$k$-characterization. In the case of affine-invariant linear families some special characterizations
are known to lead to $k$-testability. We describe these special characterizations next.

\begin{definition}[$k$-single-orbit characterization]\label{defn:single orbit}
Let $C=\big(\overline \alpha, \bset{\overline \lambda_i}_{i=1}^{r}\big)$ be a $k$-constraint
on $\{\F_{q}^n \to \F_q\}$.
The \emph{orbit} of $C$ under the set of affine-transformations is the set of $k$-constraints
\[\set{T \circ C}_T = \set{ \big( \big(T(\alpha_1),\ldots,T(\alpha_k)\big), \bset{\overline \lambda_i}_{i=1}^{r} \big)\; \; \bigg| \; \; T:\F_{q}^n \to \F_q^n \; \text{is an affine-transformation}}.\]
We say that $C$ is a \emph{$k$-single-orbit characterization} of $\cF$
if the orbit of $C$ forms a $k$-characterization of $\mathcal{F}$.
\end{definition}

The following theorem, due to Kaufman and Sudan~\cite{KS08}, says that $k$-single-orbit characterization implies local testability.

\begin{theorem}[Single-orbit characterization implies local testability, {\protect~\cite[Theorem 2.9]{KS08-ECCC}}]\label{thm:KS}
Let $\mathcal{F} \subseteq \set{\F_{q}^n \to \F_q}$ be an affine-invariant linear family. If
$\cF$ has a k-single-orbit characterization, then $\cF$ has a
$(k,\epsilon)$-tester for 
$\epsilon = \min\set{1/2, {1}/{((2k+1)(k-1))}}$.

\end{theorem}

In view of the above theorem, it suffices to find a single-orbit
characterization of $\mrm[n,d,q]$ to test it. The following theorem,
which we prove in the rest of this paper, thus immediately implies
\expref{Theorem}{thm:main-imp}.

\begin{theorem}\label{thm:main}
Let $q = p^s$ for prime $p$, and let $n$, $d$ be arbitrary positive
integers.
Then the Reed-Muller code $\mrm[n,d,q]$ has a $k$-single-orbit characterization for
 \[k \leq c_q \cdot  \big(2^{p-1}+p-1\big)^{(d+1)/(q(p-1))} \cdot q^{(d+1)/q}\,,\]
where $c_q \leq 3q^4$.
\end{theorem}

\subsection{Constraints vs.\ monomials}

One of the main simplifications derived from the study
of affine-invariant linear properties is that it suffices to
analyze the performance of constraints on ``monomials''
as opposed to general polynomials.
This allows us to rephrase our target (a single-orbit
characterization of $\mrm[n,d,q]$) in somewhat simpler
terms. Below we describe some of the essential notions,
namely the ``degree set,'' the ``border set'' and the
relationship of these to single-orbit characterizations.
This leads to a further reformulation of our main theorem
as \expref{Theorem}{thm:main2}.
Variations of most of the results and notions presented in this
section appeared in previous works~\cite{KS08,GKS09,BS11,BGMSS11}.
In all the above works, with the exception of~\cite{KS08},
the notions were specialized to the case of univariate functions
mapping $\F_{q^n}$ to $\F_q$
that are invariant over the set of affine-transformations over $\F_{q^n}$.
In this work we focus on these notions in the context of
affine-invariant linear properties over the domain $\F_q^n$.

Let $\calf \subseteq \{\F_{q}^n \to \F_q\}$ be a linear affine-invariant family of
functions. Note that every member of $\{\F_{q}^n \to \F_q\}$ can be written uniquely as
a polynomial in $\F_q[x_1,x_2,\ldots,x_n]$ of degree at most $q-1$ in each variable.
For a monomial  $\prod_{i=1}^nx_i^{d_i}$ over $n$ variables,  we define its degree to be the vector $\overline d = (d_1,d_2, \ldots, d_n)$ and we define its \emph{total degree} to be $\sum_{i=1}^n d_i$.
For a function
$f:\F_{q}^n \to \F_q$ we denote its \emph{support}, denoted $\supp(f)$, to
be the set of degrees in the support of the associated polynomial. I.\,e.,
$\supp(f) = \{ \overline d \in \{0,\ldots,q-1\}^n\;|\;c_{\overline d} \ne 0\}$ where $f(x) = \sum_{\overline d} c_{\overline d} x^{\overline d}$.
The \emph{degree set} $\Deg(\cF)$ of $\calf$ is simply the union of the supports of the functions
in $\calf$, \ie, $\Deg(\calf) = \cup_{f \in \calf} \supp(f)$.

While the degree set of the Reed-Muller codes are natural to study, they
are also natural in more general contexts.
The following lemma from~\cite{KS08} says that every affine-invariant linear property from $\F_q^n$ to $\F_q$ is uniquely determined by its degree set.

\begin{lemma}[Monomial extraction lemma, {\protect~\cite[Lemma 4.2]{KS08-ECCC}}]\label{lem:monomial-extraction}
Let $\mathcal{F} \subseteq \set{\F_{q}^n \to \F_q}$ be an affine-invariant linear property. Then $\cF$ has a monomial basis, that is, $\cF$ is the set of
all polynomials supported on monomials of the form $x^{\overline d}$ where $\overline d  \in \Deg(\cF)$.\footnote{Our language is somewhat different from that of~\cite{KS08-ECCC}. After
translation, their lemma says that all monomials $x^{\overline d}$
are contained in $\cF$. The other direction saying $\cF$ is contained
in the span of such monomials is immediate from the definition of
$\Deg(\cF)$.}
\end{lemma}

One main structural feature of the degree sets of affine-invariant linear properties is that
they are \emph{$p$-shadow-closed}. Before giving the definition of a shadow-closed set of degrees we need to introduce a bit of notation.
For a pair of integers $a,b$ let
$a= \sum_j a_j p^j$, $b= \sum_j b_jp^j$ be their base-$p$ representation, respectively. We say that $b$ is in the \emph{$p$-shadow} of $a$, and denote this $b \leq_p a$, if $b_j \leq a_j$ for all $j$.
For a pair of integer vectors $ \overline d=(d_1,d_2,\ldots,d_n)$, $\overline e=(e_1,e_2,\ldots,e_n)$ we say that
$\overline e \leq_p  \overline d$ if
$e_i \leq_p d_i$ for every $i$.

\begin{definition}[Shadow-closed set of degrees]\label{defn:shadow-closed}
For a vector of integers $\overline d=(d_1,d_2,\ldots,d_n)$ of length $n$, the \emph{$p$-shadow} of $\overline d$ is the set
$\Shadow_p(\overline d) = \{\overline e=(e_1,e_2,\ldots,e_n) \mid \overline e \leq_p \overline d\}$.
For a subset $S$ of integer vectors  of length $n$ we let
$\Shadow_p(S) = \bigcup_{\overline d \in S} \Shadow_p(\overline d)$.
Finally, we say that $S$ is \emph{$p$-Shadow-closed} if $\Shadow_p(S) = S$.
\end{definition}

The following lemma from~\cite{KS08-ECCC} says that the degree set of every affine-invariant linear property over $\F_q^n$ is $p$-shadow-closed.

\begin{lemma}[Monomial spread lemma, {\protect~\cite[Lemma 4.6]{KS08-ECCC}}]\label{lem:monomial-spread}
%
Let $\mathcal{F} \subseteq \set{\F_{q}^n \to \F_q}$ be an affine-invariant linear property. Then $\Deg(\cF)$ is $p$-shadow-closed.
\end{lemma}

\begin{remark}
Note that Lemma 4.6 in~\cite{KS08-ECCC} actually shows that $\Deg(\cF)$ is closed under an even stronger notion of $p$-shadow in which a pair of integer vectors $ \overline d=(d_1,d_2,\ldots,d_n)$, $\overline e=(e_1,e_2,\ldots,e_n)$ is said to satisfy
$\overline e \leq_p  \overline d$ if $\sum_{i=1}^n e_{i,j} \leq \sum_{i=1}^n d_{i,j}$ for all $j$, where $e_i = \sum_{j} e_{i,j} p^j$,  $d_i = \sum_{j} d_{i,j} p^j$ are the base-$p$ representations
of the integers $d_i, e_i$ respectively. It can be verified that if $\overline d, \overline e$ satisfy $\overline e \leq_p \overline d$ according to our definition then they also satisfy  $\overline e \leq_p \overline d$   according to the definition of~\cite{KS08-ECCC}. The converse, however, is false.
\end{remark}

A central element used in the proof of the above lemma and
other aspects in the study of affine-invariant linear properties is
Lucas's theorem, which we also need.

\begin{theorem}[Lucas's Theorem]\label{thm:luca}
The %
binomial
coefficient $\binom {n} {i} $ is non-zero mod $p$ if and only if $i \leq_p n$.
\end{theorem}

The fact that the degree set of a linear affine-invariant family is $p$-shadow-closed 
motivates the notion of a ``border'' set, the set of minimal elements (under $\leq_p$) that are not in $\Deg(\cF)$.
\begin{definition}[Border]\label{defn:border}
For an affine-invariant linear family $\calf \subseteq \set{\F_q^n \to \F_q}$, its \emph{border set}, denoted $\Border(\calf)$,
is the set \[\Border(\calf) = \{\overline e \in \{0,\ldots,q-1\}^n \;|\; \overline e \notin \Deg(\calf) \text{~but~}
\forall \overline e' \leq_p \overline e, \overline e'\ne \overline e, \overline e' \in \Deg(\calf)\}\,.\]
\end{definition}

The relationship between the degree set and the border set of an affine-invariant linear family and single-orbit characterization is given by the following lemma.
This lemma  says that for an affine-invariant linear family, in order to establish $k$-single-orbit characterization it suffices to exhibit a $k$-constraint whose orbit accepts
all monomials of the form $x^{\overline d}$ for
 $\overline d \in \Deg(\cF)$ and  rejects all monomials of the form $x^{\overline b}$ for
  $\overline b \in \Border(\cF)$. It is similar in spirit to Lemma 3.2 of~\cite{BGMSS11} which shows that a similar result holds for affine-invariant linear properties over $\F_{q^n}$.

\begin{lemma}\label{lem:border-single-orbit}
Let $\mathcal{F} \subseteq \{\F_q^n \to \F_q\}$ be an
affine-invariant linear property and let $C$ be a constraint.
Then $C$ is a single-orbit characterization of $\calf$ if the orbit of
$C$ accepts every monomial $x^{\overline d}$ for $\overline d \in \Deg(\calf)$ and
rejects every monomial $x^{\overline b}$ for $\overline b \in \Border(\calf)$.
\end{lemma}

\begin{proof}
We need to show that for every affine-transformation $T:\F_q^n \to \F_q^n$ the constraint $T \circ C$ accepts all functions $f \in \cF$,
while for every $f \notin \cF$ there exists an affine-transformation $T$ such that $T \circ C$ rejects $f$.

Since the set of monomials $x^{\overline d}$ for $\overline d \in \Deg(\cF)$ forms a basis for $\cF$, clearly we have that $C$ accepts all functions $f \in \cF$. The fact that $\cF$ is affine-invariant implies in turn that for every affine-transformation $T$ the constraint $T \circ C$ also accepts all functions $f \in \cF$.

It remains to show that for every $f \notin \cF$ there exists an affine-transformation $T: \F_q^n \to \F_q^n$ such that $T \circ C$ rejects $f$.
Suppose in contrary that there exists a function $f \notin \cF$ such that the orbit of $C$ accepts $f$, and
let \[{\widetilde \calf} = \{f' : \F_q^n \to \F_q \;|\;
T \circ C \text{~accepts~} f' \text{~for~every~affine-transformation~}T\}\,. \]
Note that $\widetilde \calf$ is a linear affine-invariant property, and that our assumption on $f$ implies that $f \in {\widetilde \cF}$.

Since $f \notin \cF$, and since $\Deg(\cF)$ forms a basis for $\cF$, there exists a monomial $x^{\overline e}$ in the support of $f$ such that $\overline e \notin \Deg(\cF)$. The 
%
Monomial Extraction Lemma %
(\expref{Lemma}{lem:monomial-extraction}) then implies that $x^{\overline e}$ is also contained in $\widetilde \cF$. Let $\overline b$ be a minimal degree (with respect to $\leq_p$) such that $\overline b \leq_p \overline e$ and $\overline b \notin \Deg(\cF)$. Then from the definition of the border we have that $\overline b \in \Border(\cF)$. Furthermore, since $\widetilde \cF$ is linear and affine-invariant and $\overline e \in \Deg({\widetilde \cF})$,
the monomial spread lemma (\expref{Lemma}{lem:monomial-spread}) implies that $\overline b \in \Deg({\widetilde \cF})$ and in particular $x^{\overline b} \in \widetilde \cF$.
But this implies in turn that  $x^{\overline b}$ is accepted by the orbit of $C$, a contradiction to our assumption that all degrees in the 
%
border of
$\cF$ are rejected by the orbit of $C$.
\end{proof}

In order to describe the border of the Reed-Muller family we shall use the following definition.
\begin{definition}\label{defn:bi}
For %
an
integer $d$, let $d_0,d_1,\ldots$ be its 
%
base-$p$ expansion, \ie, %
the $d_j$ satisfy   %
$0 \leq d_j < p$ and
$d = \sum_{j=0}^{\infty} d_j p^j$.
Let $b_i(d) = p^i + \sum_{j=i}^\infty d_j p^j$.
\end{definition}

Note that $b_i(d) > d$ for every $i$ and
conversely, for every integer $e > d$ there exists an index $i$ such
that $b_i(d) \leq_p e$.
%
The $b_i(d)$      %
are useful in describing
the border monomials of the Reed-Muller family, as formalized below.

\begin{proposition}
\label{prop:rmborder}
For every $n,d,q$, where $q=p^s$ for a prime $p$, we have
\begin{align*}
  \Deg\big(\mrm[n,d,q]\big) & =  \set{\overline d=(d_1,\ldots,d_n) \in \{0,\ldots,q-1\}^n \; \bigg| \; \sum_{j=1}^n d_j \leq d
} \quad\text{and}\\
\Border\big(\mrm[n,d,q]\big) & \subseteq 
\set{\overline e=(e_1,\ldots,e_n) \in \{0,\ldots,q-1\}^n \; \bigg| \;  \sum_{j=1}^n e_j = b_i(d) \; \text{for some} \; 0 \leq i \leq s   }\,.
\end{align*}
\end{proposition}

\begin{proof}
The fact that the degree set contains all ${\overline d}$
with $\sum_{j=1}^n d_j \leq d$ is immediate from the definitions.
Now consider ${\overline f}$ such that $\sum_{j=1}^n f_j > d$.
To verify the correctness of the border,
we wish to show that there exists ${\overline e} \leq_p
{\overline f}$ and $0 \leq i \leq s$ such that $\sum_{j=1}^n e_j =
b_i(d)$. Let $\ell$ be the least index such that $\sum_{j=1}^\ell
f_j = f > d$. Note that $f \leq d + q - 1$, since $f_\ell \leq
q - 1$. Now let $f^{(0)},f^{(1)},\ldots$ denote the
base-$p$ expansion of $f$ and let
$d^{(0)},d^{(1)},\ldots$ denote the base-$p$ expansion of $d$.
Since $f > d$, there must exist a largest index $i$ such that
$f^{(i)} > d^{(i)}$ and for all $j > i$, $f^{(j)} = d^{(j)}$.
For this choice of $i$, note that $b_i(d) \leq _p f$
and
one can reduce %
each $f_j$ to the corresponding $e_j$ %
so that
$\overline e \leq_p \overline f$ and $\sum_{j=1}^n e_j = b_i(d)$.

It remains to be shown that $i \leq s$. For this part
note that if $b_i(d) > b_{i-1}(d)$ then $b_i(d) \geq b_{i-1}(d) + p^{i-1}$.
Thus for all $i > s$, we have either $b_i(d) = b_s(d)$
or $b_i(d) \geq b_s(d) + p^s > d + q$. But since $f \leq d + q-1$
it follows that we never need to use $i > s$.
\end{proof}

Combining \expref{Lemma}{lem:border-single-orbit}
and \expref{Proposition}{prop:rmborder}
we have that \expref{Theorem}{thm:main}
follows immediately from \expref{Theorem}{thm:main2} below.

\begin{theorem}\label{thm:main2}
Let $q=p^s$ for a prime $p$.
Then there exists a $k$-constraint $C$ on $\set{\F_q^n \to \F_q}$ whose orbit
accepts all monomials
of total degree at most $d$
and rejects all
monomials of total degree $b_i(d)$ for $0 \leq i \leq s$,
for \[k \leq 3q^4 \cdot \big(2^{p-1}+p-1\big)^{(d+1)/(q(p-1))} \cdot q^{(d+1)/q}\,.\]
\end{theorem}

The rest of this paper will be devoted to proving \expref{Theorem}{thm:main2}.

